% problem-000164 6 非水电极电势热力学
\section*{6. 非水电极电势热力学}
电镀铜是电⼦⼯业中的关键技术之⼀。电镀铜的基本电化学过程如下：

\subsection*{(1)}
Cu2+ + e− → Cu+（慢） (2) $\mathrm { C u ^ { + } + e ^ { - }  C u \ ( \mathbb { H } ) }$

在电解体系中引⼊各种有机配体或者添加剂可以改进电解效率和镀层品质。末端含有磺酸基团的短碳链的⼆硫化物或硫醇（分⼦式⻅下表），如SPS（⼆硫⼆丙烷磺酸）和MPS（3-硫基-1-丙烷磺酸），常⽤做镀铜⼯艺的加速剂。有研究认为SPS或MPS加速电镀铜沉积的原因是因为它们能促进Cu+还原为Cu的反应。

（图：）

鉴于难以采⽤电化学⽅法直接测定该电对的电极电势，研究⼈员设计了巧妙的⽅法，通过1HNMR谱的定量分析（谱峰⾯积正⽐于氢原⼦的数⽬），间接求算电极电势。该⽅法具体实验如下：298K下，在氮⽓保护下将SPS和ME（巯基⼄醇）在密封容器中混合，⼆者起始浓度分别为10.0mM和20.0mM。待反应达平衡后，测试混合体系的1H-NMR，发现MPS和SPS的β位氢原⼦的谱峰（分别在 2.02 ppm 和 2.15 ppm 处）⾯积⽐为 1.00:1.00。ME 氧化产物为 DE，已知 $1 E ^ { \ominus } \mathrm { ( D E / M E ) }$ $= 0 . 1 5 3 ~ \mathrm { V _ { c } }$

\subsection*{6-1}
写出SPS和ME反应的⽅程式。计算该反应的平衡常数。（提⽰：可以采⽤简写符号。）

\subsection*{6-2}
计算电极反应 $\mathrm { S P S } + 2 \mathrm { H ^ { + } } + 2 \mathrm { e ^ { - } }  2 \mathrm { M P S }$ 的标准电极电势。

\subsection*{6-3}
实验中观察到 $\mathrm { C u ^ { 2 + } }$ +可以将MPS氧化为SPS并形成配合物Cu(MPS)（反应1）⽽吸附在镀件表⾯，Cu(MPS)获得电⼦在镀件表⾯发⽣电沉积并放出MPS（反应2），写出反应⽅程式。
